% Preambulo comun de las auditorias del informe E1 (revision_e1/)
\documentclass[11pt,a4paper]{article}
\usepackage[utf8]{inputenc}
\usepackage[T1]{fontenc}
\usepackage[spanish,es-noshorthands]{babel}
\usepackage{lmodern}
\usepackage{microtype}
% Sin patrones de guionado en espanol en esta instalacion (ver A10): se
% desactiva el guionado automatico para evitar cortes con reglas del ingles.
\hyphenpenalty=10000
\exhyphenpenalty=10000
\usepackage[a4paper,margin=2.2cm]{geometry}
\usepackage{booktabs,array,longtable}
\usepackage{graphicx,caption,subcaption}
\usepackage{xcolor}
\usepackage{enumitem}
\usepackage{fancyhdr}
\usepackage{amsmath}
\usepackage{tikz}
\usetikzlibrary{arrows.meta,positioning}
\usepackage{natbib}
\usepackage{xurl}
\usepackage[hidelinks]{hyperref}
\urlstyle{tt}
\newcommand{\archivo}[1]{\url{#1}}
\newcolumntype{L}[1]{>{\raggedright\arraybackslash}p{#1}}
\definecolor{alta}{RGB}{180,30,30}
\definecolor{media}{RGB}{200,120,0}
\definecolor{baja}{RGB}{60,110,160}
\definecolor{cajafondo}{RGB}{245,247,250}
\newcommand{\sevA}{\textcolor{alta}{\textbf{Alta}}}
\newcommand{\sevM}{\textcolor{media}{\textbf{Media}}}
\newcommand{\sevB}{\textcolor{baja}{\textbf{Baja}}}
\setlength{\parskip}{0.35em}
\setlength{\emergencystretch}{3em}
\setlength{\parindent}{0pt}
\pagestyle{fancy}
\fancyhf{}
\fancyhead[L]{\small Auditoría del Informe del Entregable~1 -- \auditoriaid}
\fancyhead[R]{\small \thepage}
\renewcommand{\headrulewidth}{0.4pt}
% Entorno de hallazgos: ID | Severidad | Ubicacion | Hallazgo y evidencia | Correccion
\newenvironment{hallazgos}{%
\small
\begin{longtable}{@{}L{1.15cm}L{1.25cm}L{1.9cm}L{6.0cm}L{\dimexpr\textwidth-10.3cm-10.6\tabcolsep\relax}@{}}
\toprule
\textbf{ID} & \textbf{Sev.} & \textbf{Ubicación} & \textbf{Hallazgo y evidencia} & \textbf{Corrección aplicada} \\
\midrule\endfirsthead
\toprule
\textbf{ID} & \textbf{Sev.} & \textbf{Ubicación} & \textbf{Hallazgo y evidencia} & \textbf{Corrección aplicada} \\
\midrule\endhead
\bottomrule\endfoot
}{\end{longtable}}
% Macros compartidas con el informe revisado (las secciones propuestas las usan)
% Macros compartidas entre el informe revisado y las auditorias.
% Cifras verificadas contra los archivos del repositorio (ver A01/A04/A07):
\newcommand{\nUniverso}{102}   % source_id CMIP6 con tos/Omon en ESGF (00b, model_availability_priority.csv)
\newcommand{\nSeed}{41}        % publican los cuatro experimentos (config/models_seed_cmip6.csv)
\newcommand{\nSel}{40}         % completos + QC PASS (models_inventory_final.csv)
\newcommand{\nComb}{160}       % combinaciones modelo-experimento PASS (qc_report.csv)
\newcommand{\nNoEnc}{44}       % config/models_no_encontrado_44.csv
\newcommand{\nParcial}{18}     % 102 - 40 - 44
% Fecha de la portada: fija (no \today), editable en un solo lugar.
\newcommand{\fechaentrega}{28 de septiembre de 2026}
% Estado de codigo que documenta el E1 (antes de los cambios del E2)
\newcommand{\commitEone}{\texttt{8ee6184}}

% En las auditorias el texto propuesto se muestra fuera del informe: las
% referencias cruzadas se imprimen con el nombre de la seccion destino.
\makeatletter
\def\etq#1#2{\expandafter\def\csname etq@#1\endcsname{#2}}
\etq{sec:resumen}{Resumen}
\etq{sec:alcance}{Alcance}
\etq{sec:marco}{Revisión científica}
\etq{sec:antecedentes}{Antecedentes}
\etq{sec:referencias}{Literatura sobre TOE}
\etq{tab:referencias}{bibliografía TOE}
\etq{sec:similitudes}{Niño~3.4 frente a Niño~1+2}
\etq{sec:toe_metodos}{Elección de los métodos de TOE}
\etq{tab:metodos_toe}{métodos TOE}
\etq{eq:szabo_signal}{ecuación señal M1}
\etq{eq:szabo_V}{ecuación V(t)}
\etq{eq:szabo_T}{ecuación T(t)}
\etq{eq:szabo_toe}{ecuación TOE M1}
\etq{eq:tod}{ecuación ToD}
\etq{sec:datos}{Datos}
\etq{fig:cmip6_contribution}{contribución institucional}
\etq{sec:tabla_modelos}{Modelos seleccionados}
\etq{tab:modelos}{modelos}
\etq{sec:desviacion}{Selección de modelos}
\etq{sec:dominios_periodos}{Dominios y periodos}
\etq{fig:region_nino}{regiones Niño}
\etq{sec:arquitectura}{Arquitectura del pipeline}
\etq{fig:flujograma}{Flujograma}
\etq{sec:scripts}{Descripción de los pasos}
\etq{sec:paso00b}{Paso 00b}
\etq{sec:paso01}{Paso 01}
\etq{sec:paso02}{Paso 02}
\etq{sec:paso04}{Paso 04}
\etq{sec:paso05}{Paso 05}
\etq{sec:paso06}{Paso 06}
\etq{sec:paso07}{Paso 07}
\etq{sec:reporte}{Reporte de estado}
\etq{sec:resultados}{Resultados}
\etq{fig:qc}{matriz de estado}
\etq{tab:cuentas}{conteo de modelos}
\etq{fig:mapas2d}{mapas 2D}
\etq{fig:series}{series}
\etq{fig:boxplot}{boxplots}
\etq{sec:roadmap}{Próximos pasos}
\etq{tab:roadmap}{próximos pasos}
\etq{sec:conclusiones}{Conclusiones}
\etq{sec:recomendaciones}{Recomendaciones}
\makeatother

\makeatletter
\AtBeginDocument{\renewcommand{\ref}[1]{\@ifundefined{etq@#1}{\textit{[#1]}}{\textit{\csname etq@#1\endcsname}}}}
\makeatother
% Texto propuesto: secciones degradadas un nivel y en caja
\newenvironment{propuesto}{%
  \par\medskip\noindent\textcolor{baja}{\rule{\textwidth}{0.8pt}}\par
  \noindent{\small\textbf{Inicio del texto propuesto para el informe revisado}}\par\medskip
  \begingroup
  \let\section\subsection\let\subsection\subsubsection\let\subsubsection\paragraph
}{%
  \endgroup
  \par\noindent{\small\textbf{Fin del texto propuesto}}\par
  \noindent\textcolor{baja}{\rule{\textwidth}{0.8pt}}\par\medskip}
% Recuadro de actualizacion posterior (decisiones del autor)
\newcommand{\actualizacion}[1]{\par\medskip\noindent\fcolorbox{media}{media!8}{\parbox{\dimexpr\linewidth-2\fboxsep-2\fboxrule\relax}{\small\textbf{Actualización del 29-09-2026 (decisiones del autor).} #1}}\par\medskip}
